\section{Arquitetura da reorganização do projeto}
\label{sec:estrutura}

\subsection{Visão geral da hierarquia de classes adotadas}

Para substituir a implementação procedural do jogo, o projeto foi reorganizado
em uma estrutura baseada em classes e interfaces, separando responsabilidades e
agrupando comportamentos semelhantes. A organização final dos pacotes ficou da
seguinte forma:

\begin{itemize}
    \item \texttt{game.core -} reúne interfaces utilizadas por diferentes partes do jogo (\texttt{Drawable}, \texttt{Updatable}, \texttt{Damageable}).
    \item \texttt{game.core.collision -} concentra a lógica relacionada ao sistema de colisão (\texttt{Collidable} e \texttt{CollisionManager}).
    \item \texttt{game.core.entities -} contém as entidades que representam os elementos do jogo, com os subpacotes \texttt{enemies} (inimigos comuns e chefes) e \texttt{powerUps} (os dois tipos de power-up).
    \item \texttt{game.config -} reúne a leitura dos arquivos de configuração, com os subpacotes \texttt{spawn} (representação unificada das linhas de fase) e \texttt{phase} (controle de quando cada entidade deve surgir).
    \item \texttt{game.ui -} reúne os elementos de interface visual do jogo, como o \texttt{HUD} e o fundo de estrelas.
\end{itemize}

A hierarquia de classes concretas ficou assim:

\begin{center}
    WOW IMAGEM DE UML FODÃO
\end{center}

A hierarquia foi construída de forma que cada nível acrescente apenas as
características específicas daquele conjunto de entidades (princípio classico de OO). Assim, informações
comuns, como posição, movimentação ou vida, são definidas apenas um vez e
reaproveitadas pelas subclasses, reduzindo duplicação de código (consequentemente reduzindo chances de bugs) e facilitando a
manutenção da implementação.

\subsection{Interfaces em \texttt{game.core}}

As interfaces foram utilizadas para representar capacidades independentes das
entidades do jogo. Em vez de concentrar todos os comportamentos em uma única
classe base, cada interface define apenas uma responsabilidade específica (similar a um contrato),
permitindo que cada classe implemente somente os comportamentos que são realmente necessários.

\begin{itemize}
    \item \texttt{Drawable}. define a operação responsável pelo desenho da entidade na tela.
    \item \texttt{Updatable}. representa objetos cujo estado precisa ser atualizado a cada ciclo do jogo.
    \item \texttt{Damageable}. reúne as operações relacionadas ao sistema de dano e vida, sendo implementada apenas pelas entidades que podem sofrer dano.
    \item \texttt{Collidable}. define as informações necessárias para o sistema de colisão e disponibiliza um método padrão para verificar a sobreposição entre dois objetos.
\end{itemize}

Essa separação torna a arquitetura mais flexível, pois evita que uma classe
precise implementar comportamentos que não fazem sentido para ela. Por exemplo,
os elementos do \texttt{game.ui}, como o fundo de estrelas, implementam apenas
\texttt{Drawable} e \texttt{Updatable}, sem participar do sistema de colisão
nem possuir pontos de vida, já que essas capacidades não fazem sentido para um
elemento puramente visual.

\begin{lstlisting}[caption={contrato e comprtamento padrão da interface \texttt{Collidable} (Collidable.java)}]
public interface Collidable {
    double getX();
    double getY();
    double getWidth();
    double getHeight();

    default boolean hasCollidedWith(Collidable other){
        double halfW1 = getWidth() / 2.0;
        double halfH1 = getHeight() / 2.0;
        double halfW2 = other.getWidth() / 2.0;
        double halfH2 = other.getHeight() / 2.0;

        return Math.abs(getX() - other.getX()) < (halfW1 + halfW2)
            && Math.abs(getY() - other.getY()) < (halfH1 + halfH2);
    }
}
\end{lstlisting}

Por ser um método \texttt{default}, qualquer classe que implemente
\texttt{Collidable} já ganha automaticamente a lógica de detecção de
sobreposição entre retângulos, sem precisar reimplementar.

\subsection{Por que \texttt{GameObject} e \texttt{Entity} são classes separadas}

A classe \texttt{GameObject} reúne apenas as informações presentes em qualquer
objeto exibido na tela, como posição, estado de atividade e a verificação de
quando um objeto saiu dos limites da tela:

\begin{lstlisting}[caption={Responsabilidades mínimas comuns a qualquer objeto do jogo (GameObject.java)}]
public abstract class GameObject implements Drawable, Updatable {
    protected double x;
    protected double y;
    protected boolean active;
    ...
    protected boolean isOutOfBounds(){
        return x < - OUT_OF_BOUNDS_MARGIN
            || x > GameLib.WIDTH + OUT_OF_BOUNDS_MARGIN
            || y < - OUT_OF_BOUNDS_MARGIN
            || y > GameLib.HEIGHT + OUT_OF_BOUNDS_MARGIN;
    }
}
\end{lstlisting}

A classe \texttt{Entity} estende essa estrutura (ou seja \texttt{Entity} é um \texttt{GameObject}) adicionando atributos
relacionados à movimentação e às colisões, como velocidade e dimensões, além de
implementar \texttt{Collidable}. Parte da estrutura de Entity a seguir:

\begin{lstlisting}[caption={\texttt{Entity} reparoveita isOutOfBounds() da sua superclasse (Entity.java)}]
public abstract class Entity extends GameObject implements Collidable {
    protected double vx, vy;
    protected double width, height;
    ...
    @Override
    public void update(long delta){
        move(delta);

        if (isOutOfBounds()){
            deadactivate();
        }
    }
}
\end{lstlisting}

Repare que \texttt{update} cuida para que qualquer entidade
que saia da área de jogo (um projétil ou um inimigo comum, por exemplo) seja
desativada automaticamente por esse comportamento herdado (\texttt{isOutOfBounds}), sem que cada
subclasse precise implementar essa verificação por conta própria.

Essa divisão permite que diferentes tipos de entidade compartilhem toda a lógica
básica de movimentação sem carregar responsabilidades que não lhes pertencem.
Por exemplo, projéteis e power-ups herdam o comportamento de movimento definido
em \texttt{Entity}, mas não possuem atributos relacionados à vida ou dano, que
só passam a existir na classe \texttt{Ship}, utilizada pelo jogador e pelos
inimigos.

Com essa organização, a lógica comum de atualização permanece concentrada em um
único ponto da hierarquia e é automaticamente reutilizada pelas subclasses.

\subsection{Um único algoritmo de colisão para todos os pares}

A interface \texttt{Collidable} permitiu unificar a lógica de colisão em um
único algoritmo capaz de operar sobre qualquer objeto que implemente essa
interface. Dessa forma, o mesmo mecanismo pode ser utilizado para verificar
colisões entre jogadores, inimigos, projéteis e power-ups. A classe
\texttt{CollisionManager} concentra todas essas verificações:

\begin{lstlisting}[caption={Alguns métodos de \texttt{CollisionManager} operando sobre listas de classes \texttt{Collidable} (CollisionManager.java)}]
public void checkCollisions(Player player, List<Enemy> enemies, List<Projectile> projectiles){
    checkProjectilesAgainstEnemies(projectiles, enemies);
    checkProjectilesAgainstPlayer(projectiles, player);
    checkEnemiesAgainstPlayer(enemies, player);
}

public void checkPowerUpPickup(Player player, List<PowerUp> powerUps){
    for (PowerUp p : powerUps){
        if (p.isActive() && p.hasCollidedWith(player)) {
            p.collectBy(player);
        }
    }
}
\end{lstlisting}

Essa abordagem elimina a necessidade de criar uma implementação específica para
cada combinação de entidades, como ocorria no código original, no qual cinco
laços praticamente idênticos verificavam colisões entre pares diferentes de
vetores. Sempre que um novo objeto participa do
sistema de colisão, basta que ele implemente \texttt{Collidable} para que passe
a utilizar o algoritmo já existente, sem exigir nenhuma alteração no
\texttt{CollisionManager}.

\subsection{Uso de Coleções Java}

Enquanto o \texttt{Main.java} original mantinha cada conjunto de entidades em
vetores de tamanho fixo, controlando manualmente quais posições estavam
ocupadas ou livres, a reescrita aplicada na nova organização
substitui essa abordagem pelo uso do arcabouço de coleções do Java. A classe
\texttt{Game} passou a armazenar inimigos, projéteis e power-ups em listas.

\begin{lstlisting}[caption={Gerenciamento de entidades}]
private final List<Enemy> enemies;
private final List<Projectile> projectiles;
private final List<PowerUp> powerUps;
...
this.enemies = new LinkedList<>();
this.projectiles = new LinkedList<>();
this.powerUps = new LinkedList<>();
\end{lstlisting}

Como as coleções do Java são estruturas dinâmicas, não existe mais um limite
pré-estabelecido para o número de entidades simultâneas em cena, nem a necessidade
de uma busca manual por um índice livre, como fazia era feito no código original. Basta adicionar um novo elemento à
coleção correspondente quando a entidade é criada.

A remoção de entidades inativas também deixou de depender de verificações
manuais de estado espalhadas pelo código, sendo concentrada em uma única
chamada por coleção a cada ciclo de atualização.

\begin{lstlisting}[caption={Remoção de entidades inativas via \texttt{removeIf} (Game.java)}]
enemies.removeIf(e -> !e.isActive());
projectiles.removeIf(p -> !p.isActive());
powerUps.removeIf(p -> !p.isActive());
\end{lstlisting}

Dessa forma, garante-se que nenhuma entidade inativa permaneça ocupando espaço
indefinidamente, assim evitando preocupações quanto ao vazamento de memória.

Vale ressaltar que o uso de coleções não se restringiu ao gerenciamento das entidades em jogo.
A classe \texttt{Phase}, responsável por controlar a aparição de inimigos,
chefes e power-ups ao longo de uma fase, também é construída sobre uma lista.

\begin{lstlisting}[caption={Utilização de lista para agendar as aparições da fase (Phase.java)}]
private final List<EntitySpec> schedule; // ordered by whenMs

public List<EntitySpec> pollDueSpawns(long elapsedMs){
    List<EntitySpec> due = new ArrayList<>();

    while (nextIndex < schedule.size() && schedule.get(nextIndex).getWhenMs() <= elapsedMs) {
        due.add(schedule.get(nextIndex));
        nextIndex++;
    }

    return due;
}
\end{lstlisting}

Nesse caso, a lista ordenada por instante de aparição permite que o método
\texttt{pollDueSpawns} percorra o cronograma da fase de forma incremental,
devolvendo apenas as entidades cujo instante de aparição já foi atingido.

